% =====================================================================
%  materiais/_comum/preambulo.tex
%  Preambulo LaTeX das saidas em PDF. Injetado por pandoc com -H.
% =====================================================================
%  MVP. NAO reproduz o visual do livro de proposito: a pauta ciano de
%  25% e a FR Cursive de ../mary-livro/config/pauta.tex ficaram fora
%  desta etapa. O que se testa aqui e o CONTEUDO e o fluxo de
%  producao, nao a diagramacao final.
% ---------------------------------------------------------------------

\usepackage[a4paper,top=18mm,bottom=16mm,left=20mm,right=20mm]{geometry}
\usepackage{xcolor}

%  Century Schoolbook, a mesma familia do livro. Custa zero (ja esta
%  instalada) e e a fonte desenhada para livro escolar infantil.
\usepackage{tgschola}

%  Sem folio: sao folhas soltas para imprimir e preencher, nao um
%  caderno paginado. Numero de pagina numa ficha de uma folha e ruido.
\pagestyle{empty}

% ---------------------------------------------------------------------
%  \linhas{n}  --  n fios para a crianca escrever a mao
% ---------------------------------------------------------------------
%  PASSO DE 8mm, e nao os 19pt (6,7mm) da pauta do livro.
%
%  Aquele passo foi calibrado para a FR Cursive COMPOSTA a 13pt, que e
%  letra de imprensa imitando mao. Aqui quem escreve e uma crianca de
%  ~9 anos com lapis: caderno de pauta escolar brasileiro trabalha
%  entre 8 e 10mm. 6,7mm seria apertado demais.
%
%  TOM: preto a 15%. O README exige legibilidade em preto e branco, e
%  15% de K sai em qualquer laser domestica. Abaixo de ~10% some.
\newlength{\LinhaPasso}\setlength{\LinhaPasso}{8mm}
\newlength{\LinhaEsp}\setlength{\LinhaEsp}{0.4pt}
\newcommand{\LinhaTom}{black!15}

%  O VAO VEM ANTES DO FIO, nao depois.
%
%  Nao e detalhe de arrumacao: quem escreve a mao poe a letra EM CIMA
%  do fio, ou seja, o espaco util de uma linha e o que esta ACIMA dela.
%  Com o vao depois, cada bloco terminava com 8mm de papel morto e a
%  ficha estourava para uma terceira pagina -- eram ~6mm perdidos
%  vezes sete blocos.
\newcommand{\linhas}[1]{%
  \par\nobreak
  \begingroup
    \setlength{\parskip}{0pt}%
    \setlength{\parindent}{0pt}%
    \count255=0
    \loop\ifnum\count255<#1
      %  \penalty10000 -- o bloco de linhas e INDIVISIVEL.
      %
      %  Espaco para ESCREVER partido no meio deixa de servir: a
      %  crianca escreve as duas primeiras linhas, vira a folha e nao
      %  consegue mais reler o que acabou de escrever. E o mesmo
      %  raciocinio do \diariovazio de ../mary-livro/config/blocos.tex.
      %  Sem isto, o bloco de 3 linhas do campo "o que eu imaginava"
      %  saia 2 numa pagina e 1 na outra.
      \ifnum\count255>0 \penalty10000 \fi
      \vspace{\dimexpr\LinhaPasso-\LinhaEsp\relax}%
      %  \nointerlineskip e OBRIGATORIO, nao arrumacao.
      %
      %  \rule...\par compoe um PARAGRAFO, e entre paragrafos o TeX
      %  insere cola de entrelinha do tamanho de \baselineskip. A 11pt
      %  isso e 13,2pt somados a CADA fio: o passo real virava 36pt em
      %  vez dos 22,7pt (8mm) pedidos, e a ficha estourava de 2 para 3
      %  paginas. \nointerlineskip zera \prevdepth e a cola nao entra.
      %
      %  Medido: bloco de 3 fios custava 124pt; passou a custar 84pt.
      \nointerlineskip
      \noindent{\color{\LinhaTom}\rule{\linewidth}{\LinhaEsp}}\par
      \advance\count255 by 1
    \repeat
  \endgroup
  \par}

% ---------------------------------------------------------------------
%  \moldura{altura}  --  retangulo vazio, para desenhar dentro
% ---------------------------------------------------------------------
%  Sem titulo e sem legenda: a ficha n. 02 pede "um espaco embaixo,
%  SEM TITULO, para desenhar". Rotular o espaco ja diria o que
%  desenhar, e a regra e que a pergunta nunca contenha a resposta.
\newcommand{\moldura}[1]{%
  \par\nobreak\vspace{2mm}%
  \noindent{\color{\LinhaTom}\framebox[\linewidth][c]{\rule{0pt}{#1}}}%
  \par}

% ---------------------------------------------------------------------
%  Titulos compactos: cada milimetro conta para a ficha caber em uma
%  folha.
% ---------------------------------------------------------------------
\usepackage{titlesec}
\titlespacing*{\section}{0pt}{2mm}{1mm}
\titlespacing*{\subsection}{0pt}{3mm}{1mm}
\titleformat{\subsection}{\normalfont\large\bfseries}{}{0pt}{}
\setlength{\parskip}{1mm}
\setlength{\parindent}{0pt}
